
\documentclass{article}
\usepackage{colortbl}
\usepackage{makecell}
\usepackage{multirow}
\usepackage{supertabular}

\begin{document}

\newcounter{utterance}

\centering \large Interaction Transcript for game `dond', experiment `coop\_de', episode 12 with deepseek{-}chat{-}v3.1{-}t1.0.
\vspace{24pt}

{ \footnotesize  \setcounter{utterance}{1}
\setlength{\tabcolsep}{0pt}
\begin{supertabular}{c@{$\;$}|p{.15\linewidth}@{}p{.15\linewidth}p{.15\linewidth}p{.15\linewidth}p{.15\linewidth}p{.15\linewidth}}
    \# & $\;$A & \multicolumn{4}{c}{Game Master} & $\;\:$B\\
    \hline

    \theutterance \stepcounter{utterance}  
    & & \multicolumn{4}{p{0.6\linewidth}}{
        \cellcolor[rgb]{0.9,0.9,0.9}{
            \makecell[{{p{\linewidth}}}]{
                \texttt{\tiny{[P1$\langle$GM]}}
                \texttt{Sie spielen ein kooperatives Verhandlungsspiel, bei dem Sie sich mit einem anderen Spieler darauf einigen müssen, wie eine Reihe von Gegenständen aufgeteilt werden soll.} \\
\\ 
\texttt{Die Regeln:} \\
\texttt{(a) Sie und der andere Spieler erhalten eine Sammlung von Gegenständen. Jeder von Ihnen erhält außerdem eine geheime Wertfunktion, die angibt, wie viel Ihnen jede Art von Gegenstand wert ist.} \\
\texttt{(b) Sie tauschen Nachrichten mit dem anderen Spieler aus, um zu vereinbaren, wer welche Gegenstände bekommt. Sie können jeweils maximal 5 Nachrichten senden oder das Spiel vorzeitig beenden, indem Sie jederzeit einen geheimen Vorschlag machen.} \\
\texttt{(c) Jeder von euch wird aufgefordert, einen geheimen Vorschlag zu machen, in dem ihr die gewünschten Gegenstände in eckigen Klammern wie folgt angibt: "{[}Vorschlag: <Nummer> <Objektname>, <Nummer> <Objektname>, <...>{]}"} \\
\texttt{(d) Wenn eure Vorschläge komplementär sind, d.h. es gibt genug Gegenstände, um beide Vorschläge zu erfüllen, erhält jeder Spieler eine Punktzahl, die sich aus der Summe der Werte für die Gegenstände ergibt, die er erhalten hat. Andernfalls erhalten Sie beide null Punkte.} \\
\texttt{(e) Ihr Ziel ist es, die Summe aus Ihrer Punktzahl und der Punktzahl des anderen Spielers zu maximieren.} \\
\\ 
\texttt{Beginnen wir.} \\
\\ 
\texttt{Die Menge der verfügbaren Gegenstände ist:} \\
\\ 
\texttt{3 Lineale, 2 Eier, 1 Batterie} \\
\\ 
\texttt{Deine geheime Wertfunktion ist:} \\
\\ 
\texttt{Lineal: 1, Ei: 2, Batterie: 3} \\
\\ 
\texttt{WICHTIG: Ihre Nachrichten werden, sofern es sich nicht um einen geheimen Vorschlag handelt, direkt an den anderen Spieler übermittelt, also fügen Sie keine Antwort auf die Regeln oder einen Text zur Ankündigung Ihrer Nachricht ein. Um einen geheimen Vorschlag zu machen, verwenden Sie das angegebene Format. Verwenden Sie keine eckigen Klammern, wenn Sie mit dem anderen Spieler kommunizieren, sonst wird dies als Ihr geheimer Vorschlag interpretiert.} \\
            }
        }
    }
    & \\ \\

    \theutterance \stepcounter{utterance}  
    & \multicolumn{4}{p{0.6\linewidth}}{
        \cellcolor[rgb]{0.9,0.9,0.9}{
            \makecell[{{p{\linewidth}}}]{
                \texttt{\tiny{[P1$\rangle$GM]}}
                \texttt{Hallo! Ich schlage vor, dass wir die Gegenstände so aufteilen, dass die Gesamtpunktzahl maximiert wird. Was sind deine Werte für die Gegenstände?} \\
            }
        }
    }
    & & \\ \\

    \theutterance \stepcounter{utterance}  
    & & & \multicolumn{2}{p{0.3\linewidth}}{
        \cellcolor[rgb]{0.9,0.9,0.9}{
            \makecell[{{p{\linewidth}}}]{
                \texttt{\tiny{[GM$|$GM]}}
                \texttt{continue} \\
            }
        }
    }
    & & \\ \\

    \theutterance \stepcounter{utterance}  
    & & & \multicolumn{2}{p{0.3\linewidth}}{
        \cellcolor[rgb]{0.9,0.9,0.9}{
            \makecell[{{p{\linewidth}}}]{
                \texttt{\tiny{[GM$|$GM]}}
                \texttt{Hallo! Ich schlage vor, dass wir die Gegenstände so aufteilen, dass die Gesamtpunktzahl maximiert wird. Was sind deine Werte für die Gegenstände?} \\
            }
        }
    }
    & & \\ \\

    \theutterance \stepcounter{utterance}  
    & & \multicolumn{4}{p{0.6\linewidth}}{
        \cellcolor[rgb]{0.9,0.9,0.9}{
            \makecell[{{p{\linewidth}}}]{
                \texttt{\tiny{[GM$\rangle$P2]}}
                \texttt{Sie spielen ein kooperatives Verhandlungsspiel, bei dem Sie sich mit einem anderen Spieler darauf einigen müssen, wie eine Reihe von Gegenständen aufgeteilt werden soll.} \\
\\ 
\texttt{Die Regeln:} \\
\texttt{(a) Sie und der andere Spieler erhalten eine Sammlung von Gegenständen. Jeder von Ihnen erhält außerdem eine geheime Wertfunktion, die angibt, wie viel Ihnen jede Art von Gegenstand wert ist.} \\
\texttt{(b) Sie tauschen Nachrichten mit dem anderen Spieler aus, um zu vereinbaren, wer welche Gegenstände bekommt. Sie können jeweils maximal 5 Nachrichten senden oder das Spiel vorzeitig beenden, indem Sie jederzeit einen geheimen Vorschlag machen.} \\
\texttt{(c) Jeder von euch wird aufgefordert, einen geheimen Vorschlag zu machen, in dem ihr die gewünschten Gegenstände in eckigen Klammern wie folgt angibt: "{[}Vorschlag: <Nummer> <Objektname>, <Nummer> <Objektname>, <...>{]}"} \\
\texttt{(d) Wenn eure Vorschläge komplementär sind, d.h. es gibt genug Gegenstände, um beide Vorschläge zu erfüllen, erhält jeder Spieler eine Punktzahl, die sich aus der Summe der Werte für die Gegenstände ergibt, die er erhalten hat. Andernfalls erhalten Sie beide null Punkte.} \\
\texttt{(e) Ihr Ziel ist es, die Summe aus Ihrer Punktzahl und der Punktzahl des anderen Spielers zu maximieren.} \\
\\ 
\texttt{Beginnen wir.} \\
\\ 
\texttt{Die Menge der verfügbaren Gegenstände ist:} \\
\\ 
\texttt{3 Lineale, 2 Eier, 1 Batterie} \\
\\ 
\texttt{Deine geheime Wertfunktion ist:} \\
\\ 
\texttt{Lineal: 3, Ei: 0, Batterie: 1} \\
\\ 
\texttt{WICHTIG: Ihre Nachrichten werden, sofern es sich nicht um einen geheimen Vorschlag handelt, direkt an den anderen Spieler übermittelt, also fügen Sie keine Antwort auf die Regeln oder einen Text zur Ankündigung Ihrer Nachricht ein. Um einen geheimen Vorschlag zu machen, verwenden Sie das angegebene Format. Verwenden Sie keine eckigen Klammern, wenn Sie mit dem anderen Spieler kommunizieren, sonst wird dies als Ihr geheimer Vorschlag interpretiert.} \\
\\ 
\\ 
\texttt{Hallo! Ich schlage vor, dass wir die Gegenstände so aufteilen, dass die Gesamtpunktzahl maximiert wird. Was sind deine Werte für die Gegenstände?} \\
            }
        }
    }
    & \\ \\

    \theutterance \stepcounter{utterance}  
    & & & \multicolumn{4}{p{0.6\linewidth}}{
        \cellcolor[rgb]{0.9,0.9,0.9}{
            \makecell[{{p{\linewidth}}}]{
                \texttt{\tiny{[GM$\langle$P2]}}
                \texttt{Hallo! Meine Werte sind: Lineal 5, Ei 2, Batterie 4. Wir haben 3 Lineale, 2 Eier und 1 Batterie. Wie können wir die Gesamtpunktzahl maximieren?} \\
            }
        }
    }
     \\ \\

    \theutterance \stepcounter{utterance}  
    & & & \multicolumn{2}{p{0.3\linewidth}}{
        \cellcolor[rgb]{0.9,0.9,0.9}{
            \makecell[{{p{\linewidth}}}]{
                \texttt{\tiny{[GM$|$GM]}}
                \texttt{continue} \\
            }
        }
    }
    & & \\ \\

    \theutterance \stepcounter{utterance}  
    & & & \multicolumn{2}{p{0.3\linewidth}}{
        \cellcolor[rgb]{0.9,0.9,0.9}{
            \makecell[{{p{\linewidth}}}]{
                \texttt{\tiny{[GM$|$GM]}}
                \texttt{Hallo! Meine Werte sind: Lineal 5, Ei 2, Batterie 4. Wir haben 3 Lineale, 2 Eier und 1 Batterie. Wie können wir die Gesamtpunktzahl maximieren?} \\
            }
        }
    }
    & & \\ \\

    \theutterance \stepcounter{utterance}  
    & & \multicolumn{4}{p{0.6\linewidth}}{
        \cellcolor[rgb]{0.9,0.9,0.9}{
            \makecell[{{p{\linewidth}}}]{
                \texttt{\tiny{[P1$\langle$GM]}}
                \texttt{Hallo! Meine Werte sind: Lineal 5, Ei 2, Batterie 4. Wir haben 3 Lineale, 2 Eier und 1 Batterie. Wie können wir die Gesamtpunktzahl maximieren?} \\
            }
        }
    }
    & \\ \\

    \theutterance \stepcounter{utterance}  
    & \multicolumn{4}{p{0.6\linewidth}}{
        \cellcolor[rgb]{0.9,0.9,0.9}{
            \makecell[{{p{\linewidth}}}]{
                \texttt{\tiny{[P1$\rangle$GM]}}
                \texttt{Super, danke für deine Werte! Meine Werte sind: Lineal 1, Ei 2, Batterie 3.} \\
\\ 
\texttt{Um die Gesamtpunktzahl zu maximieren, sollten die Gegenstände zu dem Spieler gehen, der den höchsten Wert dafür hat.} \\
\\ 
\texttt{*   Du bewertest Lineale höher (5 vs. 1). Du solltest also alle 3 Lineale bekommen.} \\
\texttt{*   Ich bewerte Eier höher (2 vs. 2, gleich, aber ich nehme sie gerne) und die Batterie höher (3 vs. 4, du bewertest sie etwas höher).} \\
\\ 
\texttt{Um die Gesamtpunktzahl zu maximieren, sollte die Batterie zu dir gehen (4 Punkte für dich sind besser als 3 für mich). Die Eier können wir aufteilen, da der Wert gleich ist.} \\
\\ 
\texttt{Mein Vorschlag:} \\
\texttt{*   Du: 3 Lineale + 1 Batterie (Dein Gewinn: 3*5 + 1*4 = 15 + 4 = 19)} \\
\texttt{*   Ich: 2 Eier (Mein Gewinn: 2*2 = 4)} \\
\texttt{**Gesamtpunktzahl: 23**} \\
\\ 
\texttt{Alternativ, wenn du die Eier willst:} \\
\texttt{*   Du: 3 Lineale + 2 Eier (15 + 4 = 19)} \\
\texttt{*   Ich: 1 Batterie (3)} \\
\texttt{**Gesamtpunktzahl: 22**} \\
\\ 
\texttt{Die erste Aufteilung ist also besser. Was meinst du?} \\
            }
        }
    }
    & & \\ \\

    \theutterance \stepcounter{utterance}  
    & & & \multicolumn{2}{p{0.3\linewidth}}{
        \cellcolor[rgb]{0.9,0.9,0.9}{
            \makecell[{{p{\linewidth}}}]{
                \texttt{\tiny{[GM$|$GM]}}
                \texttt{continue} \\
            }
        }
    }
    & & \\ \\

    \theutterance \stepcounter{utterance}  
    & & & \multicolumn{2}{p{0.3\linewidth}}{
        \cellcolor[rgb]{0.9,0.9,0.9}{
            \makecell[{{p{\linewidth}}}]{
                \texttt{\tiny{[GM$|$GM]}}
                \texttt{Super, danke für deine Werte! Meine Werte sind: Lineal 1, Ei 2, Batterie 3.} \\
\\ 
\texttt{Um die Gesamtpunktzahl zu maximieren, sollten die Gegenstände zu dem Spieler gehen, der den höchsten Wert dafür hat.} \\
\\ 
\texttt{*   Du bewertest Lineale höher (5 vs. 1). Du solltest also alle 3 Lineale bekommen.} \\
\texttt{*   Ich bewerte Eier höher (2 vs. 2, gleich, aber ich nehme sie gerne) und die Batterie höher (3 vs. 4, du bewertest sie etwas höher).} \\
\\ 
\texttt{Um die Gesamtpunktzahl zu maximieren, sollte die Batterie zu dir gehen (4 Punkte für dich sind besser als 3 für mich). Die Eier können wir aufteilen, da der Wert gleich ist.} \\
\\ 
\texttt{Mein Vorschlag:} \\
\texttt{*   Du: 3 Lineale + 1 Batterie (Dein Gewinn: 3*5 + 1*4 = 15 + 4 = 19)} \\
\texttt{*   Ich: 2 Eier (Mein Gewinn: 2*2 = 4)} \\
\texttt{**Gesamtpunktzahl: 23**} \\
\\ 
\texttt{Alternativ, wenn du die Eier willst:} \\
\texttt{*   Du: 3 Lineale + 2 Eier (15 + 4 = 19)} \\
\texttt{*   Ich: 1 Batterie (3)} \\
\texttt{**Gesamtpunktzahl: 22**} \\
\\ 
\texttt{Die erste Aufteilung ist also besser. Was meinst du?} \\
            }
        }
    }
    & & \\ \\

    \theutterance \stepcounter{utterance}  
    & & \multicolumn{4}{p{0.6\linewidth}}{
        \cellcolor[rgb]{0.9,0.9,0.9}{
            \makecell[{{p{\linewidth}}}]{
                \texttt{\tiny{[GM$\rangle$P2]}}
                \texttt{Super, danke für deine Werte! Meine Werte sind: Lineal 1, Ei 2, Batterie 3.} \\
\\ 
\texttt{Um die Gesamtpunktzahl zu maximieren, sollten die Gegenstände zu dem Spieler gehen, der den höchsten Wert dafür hat.} \\
\\ 
\texttt{*   Du bewertest Lineale höher (5 vs. 1). Du solltest also alle 3 Lineale bekommen.} \\
\texttt{*   Ich bewerte Eier höher (2 vs. 2, gleich, aber ich nehme sie gerne) und die Batterie höher (3 vs. 4, du bewertest sie etwas höher).} \\
\\ 
\texttt{Um die Gesamtpunktzahl zu maximieren, sollte die Batterie zu dir gehen (4 Punkte für dich sind besser als 3 für mich). Die Eier können wir aufteilen, da der Wert gleich ist.} \\
\\ 
\texttt{Mein Vorschlag:} \\
\texttt{*   Du: 3 Lineale + 1 Batterie (Dein Gewinn: 3*5 + 1*4 = 15 + 4 = 19)} \\
\texttt{*   Ich: 2 Eier (Mein Gewinn: 2*2 = 4)} \\
\texttt{**Gesamtpunktzahl: 23**} \\
\\ 
\texttt{Alternativ, wenn du die Eier willst:} \\
\texttt{*   Du: 3 Lineale + 2 Eier (15 + 4 = 19)} \\
\texttt{*   Ich: 1 Batterie (3)} \\
\texttt{**Gesamtpunktzahl: 22**} \\
\\ 
\texttt{Die erste Aufteilung ist also besser. Was meinst du?} \\
            }
        }
    }
    & \\ \\

    \theutterance \stepcounter{utterance}  
    & & & \multicolumn{4}{p{0.6\linewidth}}{
        \cellcolor[rgb]{0.9,0.9,0.9}{
            \makecell[{{p{\linewidth}}}]{
                \texttt{\tiny{[GM$\langle$P2]}}
                \texttt{Das klingt nach einer guten Analyse. Die erste Aufteilung ergibt tatsächlich die höchste Gesamtpunktzahl. Ich bin mit deinem Vorschlag einverstanden: Ich nehme die 2 Eier und du nimmst die 3 Lineale und die 1 Batterie.} \\
\\ 
\texttt{{[}Vorschlag: 2 Ei, 0 Lineal, 0 Batterie{]}} \\
            }
        }
    }
     \\ \\

    \theutterance \stepcounter{utterance}  
    & & & \multicolumn{2}{p{0.3\linewidth}}{
        \cellcolor[rgb]{0.9,0.9,0.9}{
            \makecell[{{p{\linewidth}}}]{
                \texttt{\tiny{[GM$|$GM]}}
                \texttt{proposal} \\
            }
        }
    }
    & & \\ \\

    \theutterance \stepcounter{utterance}  
    & & & \multicolumn{2}{p{0.3\linewidth}}{
        \cellcolor[rgb]{0.9,0.9,0.9}{
            \makecell[{{p{\linewidth}}}]{
                \texttt{\tiny{[GM$|$GM]}}
                [0, 2, 0]
            }
        }
    }
    & & \\ \\

    \theutterance \stepcounter{utterance}  
    & & \multicolumn{4}{p{0.6\linewidth}}{
        \cellcolor[rgb]{0.9,0.9,0.9}{
            \makecell[{{p{\linewidth}}}]{
                \texttt{\tiny{[P1$\langle$GM]}}
                \texttt{Der andere Spieler hat einen geheimen Vorschlag gemacht. Jetzt bist du an der Reihe, einen eigenen Vorschlag zu machen.} \\
            }
        }
    }
    & \\ \\

    \theutterance \stepcounter{utterance}  
    & \multicolumn{4}{p{0.6\linewidth}}{
        \cellcolor[rgb]{0.9,0.9,0.9}{
            \makecell[{{p{\linewidth}}}]{
                \texttt{\tiny{[P1$\rangle$GM]}}
                \texttt{{[}Vorschlag: 2 Eier, 1 Batterie{]}} \\
            }
        }
    }
    & & \\ \\

    \theutterance \stepcounter{utterance}  
    & & & \multicolumn{2}{p{0.3\linewidth}}{
        \cellcolor[rgb]{0.9,0.9,0.9}{
            \makecell[{{p{\linewidth}}}]{
                \texttt{\tiny{[GM$|$GM]}}
                \texttt{proposal} \\
            }
        }
    }
    & & \\ \\

    \theutterance \stepcounter{utterance}  
    & & & \multicolumn{2}{p{0.3\linewidth}}{
        \cellcolor[rgb]{0.9,0.9,0.9}{
            \makecell[{{p{\linewidth}}}]{
                \texttt{\tiny{[GM$|$GM]}}
                [0, 2, 1]
            }
        }
    }
    & & \\ \\

    \theutterance \stepcounter{utterance}  
    & & & \multicolumn{2}{p{0.3\linewidth}}{
        \cellcolor[rgb]{0.9,0.9,0.9}{
            \makecell[{{p{\linewidth}}}]{
                \texttt{\tiny{[GM$|$GM]}}
                [[0, 2, 1], [0, 2, 0]]
            }
        }
    }
    & & \\ \\

\end{supertabular}
}

\end{document}
